%\usepackage{ifluatex}

\usepackage{amssymb,amsthm}
\usepackage[fleqn]{amsmath}

\usepackage[svgnames,x11names]{xcolor}
%Requires contrast ration 4.5 for AA
%asy: write(hex(heavygreen)); to return hex codes for asy colors
%pen lgreen=rgb("00a000"); to define via hex codes

\definecolor{defBlue}{HTML}{EEEEFF}%Blue for defn background
\definecolor{mRed}{HTML}{EB0000}%Red 3.99
\definecolor{dRed}{HTML}{E50000}%Red against lime thm background
\definecolor{mMagenta}{HTML}{D000D3}%Magenta: 3.13
\definecolor{dMagenta}{HTML}{C800CB}%Magenta against lime thm background
\definecolor{mPink}{HTML}{E5007E}%DeepPink: 3.63
\definecolor{mGreen}{HTML}{008A00}%lgreen(00a000)/heavygreen/LimeGreen/Lime: 3.48/2.45/2.11/1.37 (Green 5.13)
\definecolor{dGreen}{HTML}{007E00}%Green for defn background
\definecolor{mOrange}{HTML}{AE6200}%Orange: 1.97
%No problem: Blue, BlueViolet, Brown, Crimson, DarkOrchid, Green, MediumVioletRed, Teal

\usepackage{graphicx}
\usepackage{hyperref}
\hypersetup{
	pdftitle={Math 184 Notes},
	pdfauthor={Neil Donaldson},
} 
\usepackage[breakable,theorems,skins,hooks]{tcolorbox}

\usepackage{polyglossia}
\setmainlanguage{english}
\setotherlanguages{greek}

\usepackage{microtype}
\usepackage{unicode-math}
\setmainfont{TeX Gyre Pagella}
\setsansfont{TeX Gyre Heros}
\fontspec{inconsolata}
\setmathfont{TeX Gyre Pagella Math}
\setmathfont{TeX Gyre Pagella Math}[range="222B-"2233,Scale=1.08]
\setmathfont{TeX Gyre Pagella Math}[range=\sum,Scale=1]
\everydisplay{\Umathoperatorsize\displaystyle=4.5ex}

%Hierogylphs
\newfontfamily\Hiero{NotoSansEgyptianHieroglyphs-Regular.otf}[Scale=1.1,FakeBold=1.5]
\newfontfamily\Hierobig{NotoSansEgyptianHieroglyphs-Regular.otf}[Scale=1.8,FakeBold=1.5]

%Greek
\newfontfamily\Tgrk{Noto Serif Light}[Scale=1]
\def\lowgknum{{\Tgrk ͵}}
\def\upgknum{{\Tgrk ʹ}}
\def\gkfifty{Π\hspace{-5pt}\makebox[0pt][c]{\raisebox{2.5pt}{\scalebox{0.5}{Δ}}}\hspace{5pt}}
\def\gkfivehun{Π\hspace{-5pt}\makebox[0pt][c]{\raisebox{2.5pt}{\scalebox{0.5}{Η}}}\hspace{5pt}}
\def\gkfivethou{Π\hspace{-5pt}\makebox[0pt][c]{\raisebox{2.5pt}{\scalebox{0.5}{Χ}}}\hspace{5pt}}

%Astro
\newfontfamily\astro{Arial}[Scale=1,FakeBold=1.5]
\def\mars{\text{\astro ♂}}
\def\earth{\oplus}

%lengths
\unitlength 1cm
\setlength{\textheight}{22.5cm}
\textwidth 17.5cm
\oddsidemargin -0.75cm
\evensidemargin -0.75cm
\topmargin -1.5cm
\topskip 0cm
\headheight 0.5cm
\headsep 1cm
\marginparwidth 1.2cm
\parindent0pt


\newlength\doubleind
\addtolength{\doubleind}{\leftmargini}
\addtolength{\doubleind}{\leftmarginii}

%MathNumbering
\makeatletter
\@addtoreset{equation}{section}
\renewcommand{\theequation}{\thesection.\@arabic\c@equation}
\newcommand\Nopagebreak{\@nobreaktrue\nopagebreak}
%\renewcommand{\theequation}{\@arabic\c@equation}
\makeatother


\def\exstart{\hangindent\leftmargini\textup{1.}\hspace{\labelsep}}


%Theorems
\tcbset{
	defstyle/.style={enhanced, top=3pt, bottom=3pt, colframe=black, coltitle=black, arc=5pt, boxrule=1.5pt, left*=0pt, right*=0pt, theorem style=plain, terminator sign={.\ \ \ }, fonttitle=\bfseries\upshape, fontupper=\upshape, colback=defBlue, grow sidewards by=8pt, drop fuzzy shadow},
	exstyle/.style={enhanced, breakable, beforeafter skip balanced=10pt, coltitle=black, theorem style=plain, terminator sign={.\ \ \ }, fonttitle=\bfseries\upshape, fontupper=\upshape, blanker, borderline west={4pt}{-8pt}{orange!75!white}},
	thmstyle/.style={enhanced, top=3pt, bottom=3pt, colframe=black, coltitle=black, arc=5pt, boxrule=1.5pt,left*=0pt, right*=0pt, theorem style=plain, terminator sign={.\ \ \ }, fonttitle=\bfseries\upshape, fontupper=\slshape, colback=green!12!white, grow sidewards by=8pt, drop fuzzy shadow},
	exercisestyle/.style={enhanced, breakable, beforeafter skip balanced=10pt, coltitle=black, theorem style=plain, terminator sign={.\ \ \ }, fonttitle=\bfseries\upshape, fontupper=\upshape, blanker, borderline west={4pt}{-8pt}{purple!75!white}, title={Exercises\ \ \ }},
	proofstyle/.style={enhanced, breakable, beforeafter skip balanced=10pt, blanker, borderline west={4pt}{-8pt}{green!50!white}},
	asidestyle/.style={enhanced, breakable, beforeafter skip balanced=10pt, blanker, borderline west={4pt}{-8pt}{black!50!white}},
	exercisestyle2/.style={enhanced, breakable, beforeafter skip balanced=10pt, coltitle=black, theorem style=plain, terminator sign={.\ \ \ }, fonttitle=\bfseries\upshape, fontupper=\upshape, blanker, borderline west={4pt}{-8pt}{purple!75!white}, title={Exercises\ \thesubsection.\ \ \ }},
	exercisestyle4/.style={enhanced, breakable, beforeafter skip balanced=10pt, coltitle=black, theorem style=plain, terminator sign={.\ \ \ }, fonttitle=\bfseries\upshape, fontupper=\upshape, blanker, borderline west={4pt}{-8pt}{purple!75!white}, title={Exercises\ \thesection.\ \ \ }},
	exercisestyle3/.style={enhanced, breakable, beforeafter skip balanced=10pt, coltitle=black, theorem style=plain, terminator sign={.\ \ \ }, fonttitle=\bfseries\upshape, fontupper=\upshape, blanker, borderline west={4pt}{-8pt}{purple!75!white}, title={Exercise\quad}},
	}

\tcolorboxenvironment{proof}{breakable,proofstyle}


\newenvironment{proofpic}{%\topsep\relax
  \trivlist
  \item[\hskip\labelsep
        \itshape
    Proof.]\ignorespaces
}{\endtrivlist}


\newtcbtheorem[number within=section]{defn}{Definition}{defstyle}{defn}
\newtcbtheorem[use counter from=defn]{example}{Example}{exstyle}{ex}
\newtcbtheorem[use counter from=defn]{examples}{Examples}{exstyle}{ex}
\newtcbtheorem[use counter from=defn]{thm}{Theorem}{thmstyle}{thm}
\newtcbtheorem[use counter from=defn]{lemm}{Lemma}{thmstyle}{lemm}
\newtcbtheorem[use counter from=defn]{cor}{Corollary}{thmstyle}{cor}
\newtcbtheorem[use counter from=defn]{axiom}{Axiom}{defstyle}{axiom}
\newtcbtheorem[use counter from=defn]{axioms}{Axioms}{defstyle}{axioms}
\newtcbtheorem[use counter from=defn]{conj}{Conjecture}{thmstyle}{conj}

\newtcolorbox{aside}{asidestyle}
\newtcolorbox{exercises*}{exercisestyle}
\newtcolorbox{exercises}{exercisestyle2}
\newtcolorbox{exercisessec}{exercisestyle4}
\newtcolorbox{exercise}{exercisestyle3}
                      
\newenvironment{enumeratea}{
	\begin{enumerate}
	  \renewcommand{\labelenumi}{(\alph{enumi})}}
	{\end{enumerate}}
	
